% !TeX root = ../main.tex
\chapter{پیش‌پردازش؛ ساختن نمایش درست از داده}
\source{اسلاید پیش‌پردازش، صفحات PDF ۳ تا ۵۹؛ کتاب، فصل ۲، بخش‌های ۲٫۱ تا ۲٫۴ و فصل ۶، بخش عدم توازن طبقه‌ها.}

\section{کیفیت داده و معنای نوع ویژگی}
الگوریتم، معنای واقعی یک ستون را از نام آن نمی‌فهمد. ابتدا برای هر ویژگی نوع، واحد، دامنهٔ معتبر، روش اندازه‌گیری و زمان دسترس‌پذیری را ثبت کنید. ویژگی اسمی \lr{Nominal} فقط هویت گروه را بیان می‌کند؛ ویژگی ترتیبی \lr{Ordinal} ترتیب دارد ولی فاصلهٔ رتبه‌ها الزاماً برابر نیست؛ ویژگی فاصله‌ای مبدأ قراردادی دارد، مانند دمای سلسیوس؛ ویژگی نسبتی صفر معنادار دارد، مانند طول. مجاز بودن جمع، نسبت یا فاصله به این معنا بستگی دارد.

کتاب خطای اندازه‌گیری را از خطای گردآوری جدا می‌کند. نویز مؤلفهٔ تصادفی خطاست؛ سوگیری اندازه‌گیری می‌تواند همهٔ مقدارها را در یک جهت جابه‌جا کند. تکرارپذیری بالا به‌تنهایی صحت را تضمین نمی‌کند: ترازویی که همیشه دو کیلوگرم بیشتر نشان می‌دهد، ممکن است اندازه‌گیری‌های بسیار نزدیک به هم ولی نادرست تولید کند.

دادهٔ پرت \lr{Outlier} الزاماً خطا نیست. خرید بسیار بزرگ می‌تواند ورود اشتباه یک صفر یا یک مشتری واقعی و مهم باشد. حذف خودکار تمام نقاط نادر ممکن است همان پدیده‌ای را پاک کند که می‌خواستیم کشف کنیم. ابتدا قیود قطعی دامنه، ناسازگاری واحد، سطرهای تکراری و منشأ ثبت را بررسی کنید. حدهایی مانند «دمای مجاز» فقط در زمینهٔ مشخص معنا دارند.

\section{مقدار گمشده؛ نبودن داده هم اطلاعات دارد}
میان «صفر»، «پرسیده نشده» و «قابل اعمال نیست» تفاوت وجود دارد. صفر کردن همهٔ فقدان‌ها می‌تواند فاصله‌ها، میانگین‌ها و مرز تصمیم را عوض کند. حذف سطر، جای‌گذاری میانگین یا میانه، مقدار پرتکرار، نشانهٔ فقدان و مدل پیش‌بینی مقدار گمشده، گزینه‌هایی با فرض‌های متفاوت‌اند. میانگین معمولاً واریانس را کم و تراکم مصنوعی در مرکز ایجاد می‌کند؛ میانه در برابر چند مقدار بسیار بزرگ مقاوم‌تر است، ولی ساختار چندمتغیره را بازسازی نمی‌کند.

در \lr{MCAR} فقدان مستقل از داده‌هاست؛ در \lr{MAR} پس از شرط کردن بر اطلاعات مشاهده‌شده می‌توان سازوکار فقدان را توضیح داد؛ در \lr{MNAR} خود مقدار مشاهده‌نشده نیز در احتمال فقدان نقش دارد. این دسته‌بندی توضیح می‌دهد چرا یک جای‌گذاری یکسان برای همهٔ ستون‌ها، راه‌حل عمومی نیست. یک نشانگر دودویی برای «این مقدار گمشده بود» می‌تواند بخشی از اطلاعات را حفظ کند، ولی سوگیری گردآوری را خودبه‌خود برطرف نمی‌کند.

\begin{example}{جای‌گذاری می‌تواند رابطه را وارونه کند}
در مثال اسلاید، برای ۴۰ نفر قند خون بالا ثبت شده و ۳۰ نفرشان عارضه دارند: نسبت $30/40=0.75$ است. برای ۲۰۰ نفر قند خون نامعلوم است و ۲۰ نفر عارضه دارند. اگر همهٔ مقدارهای نامعلوم را «بالا» قرار دهیم، گروه «بالا» به ۲۴۰ نفر با ۵۰ عارضه تبدیل می‌شود و نسبت $50/240\approx0.208$ خواهد شد. دادهٔ تازه‌ای دربارهٔ قند خون کشف نکرده‌ایم؛ صرفاً با فرض جای‌گذاری، معنی گروه را عوض کرده‌ایم.
\end{example}

در جای‌گذاری مبتنی بر همسایه، افراد مشابه بر اساس ویژگی‌های موجود پیدا می‌شوند و از مقدارهایشان برای برآورد استفاده می‌شود. فاصله باید با فقدان و مقیاس سازگار باشد. تعداد همسایه و وزن‌دهی، تنظیمات مدل‌اند و باید در اعتبارسنجی انتخاب شوند. میانگین یا مدل جای‌گذاری را فقط روی آموزش برازش کنید.

\section{عدم توازن طبقه‌ها و نمونه‌سازی}
اگر از هزار نمونه فقط ده نمونه مثبت باشند، مدلی که همیشه منفی بگوید دقت $0.99$ دارد، ولی هیچ مثبت را پیدا نمی‌کند. ابتدا ماتریس درهم‌ریختگی و هزینهٔ خطاها را بررسی کنید. کم‌نمونه‌گیریِ طبقهٔ بزرگ داده را کنار می‌گذارد؛ بیش‌نمونه‌گیری با تکرار، اطلاعات مستقل تازه تولید نمی‌کند و می‌تواند به حفظ نمونه‌ها بینجامد. وزن طبقه، آستانهٔ تصمیم و معیار مناسب نیز راه‌های مقابله با این مسئله‌اند.

در \lr{SMOTE}، از یک نمونهٔ اقلیت و یکی از همسایه‌های همان طبقه، نقطه‌ای در پاره‌خط میان آن‌ها می‌سازیم:
\[\mathbf x_{\mathrm{new}}=\mathbf x+\lambda(\mathbf x_{\mathrm{neighbor}}-\mathbf x),\qquad 0\le\lambda\le1.\]
در نسخهٔ متعارف، یک ضریب مشترک برای مختصات یک نقطه استفاده می‌شود؛ انتخاب ضریب مستقل برای هر مؤلفه الزاماً همان نمونه‌سازی روی پاره‌خط نیست. نمونهٔ مصنوعی می‌تواند مرز طبقه را خراب کند، به‌ویژه اگر همسایه نامناسب یا دادهٔ آمیختهٔ عددی و اسمی داشته باشیم. بازنمونه‌گیری فقط داخل دادهٔ آموزش هر دور انجام می‌شود؛ آزمون باید توزیع واقعی کاربرد را بازتاب دهد.

\section{مقیاس‌بندی و جلوگیری از نشت اطلاعات}
فاصلهٔ دو فرد با سن‌های ۲۵ و ۳۰ و درآمدهای ۷۰۰۰۰ و ۸۵۰۰۰ تقریباً فقط از اختلاف درآمد اثر می‌گیرد. استانداردسازی و مقیاس‌بندی کمینه‌ـ‌بیشینه دو نمایش متفاوت می‌سازند:
\[z_{ij}=\frac{x_{ij}-\mu_j}{\sigma_j},\qquad u_{ij}=\frac{x_{ij}-a_j}{b_j-a_j}.\]
در این روابط، $\mu_j,\sigma_j,a_j,b_j$ از آموزش به دست می‌آیند. برای ویژگی ثابت، تقسیم بر صفر باید با قاعده‌ای مشخص، مثلاً خروجی صفر و ثبت ویژگی ثابت، مدیریت شود. آزمون می‌تواند خارج از بازهٔ آموزش باشد؛ در این صورت تبدیل کمینه‌ـ‌بیشینه ممکن است مقداری خارج از $[0,1]$ بدهد و این تناقضی ندارد.

\begin{important}{برازش تبدیل هم بخشی از آموزش است}
اگر میانگین، واژگان متن، مؤلفه‌های PCA یا ویژگی‌های منتخب را با کل داده محاسبه کنیم، اطلاعات اعتبارسنجی یا آزمون وارد فرایند یادگیری شده است. ترتیب درست عبارت است از تقسیم داده، برازش تبدیل روی آموزش و اعمال همان تبدیل روی دادهٔ ندیده. در اعتبارسنجی متقابل، این کار داخل هر دور تکرار می‌شود.
\end{important}

استانداردسازی، توزیع را الزاماً نرمال نمی‌کند؛ تنها مکان و مقیاس را تغییر می‌دهد. درخت تصمیم نسبت به تبدیل یکنوای یک ویژگی رفتاری متفاوت با فاصله‌محورها دارد، زیرا می‌تواند آستانه را جابه‌جا کند، اما روش‌هایی مانند همسایه، SVM ، PCA و رگرسیون با جریمه به مقیاس حساس‌اند. برای دادهٔ پرت، مقیاس‌بندی مقاوم با میانه و دامنهٔ میان‌چارکی می‌تواند مناسب‌تر باشد.

\labref{در آزمایشگاه مقیاس، فاصله و همسایهٔ نزدیک را در حالت خام، استاندارد و کمینه‌ـ‌بیشینه مقایسه کنید. بررسی کنید چرا تبدیل محور و تغییر تعریف فاصله باید از هم متمایز باشند.}

\section{پیش‌بین کاذب و طراحی خط لوله}
ویژگی‌ای که از پاسخ مشتق شده، یا فقط پس از رخداد هدف در دسترس است، پیش‌بین معتبر زمان استفاده نیست. رتبهٔ همان دانشگاه برای پیش‌بینی «دانشگاه برتر»، نتیجهٔ نهایی آزمون برای پیش‌بینی قبولی پیش از آزمون و وضعیت پرداخت آینده برای پیش‌بینی عدم پرداخت نمونه‌های نشت‌اند. همبستگی بالا نشانهٔ بررسی است، نه حکم قطعیِ کاذب بودن؛ زمان تولید و معنای ویژگی تعیین‌کننده‌اند.

یک \lr{Pipeline} پیش‌پردازش و مدل را به واحدی تبدیل می‌کند که در هر دور اعتبارسنجی از نو برازش می‌شود. \lr{ColumnTransformer} برای اعمال تبدیل متفاوت روی ستون‌های عددی، اسمی و متنی است. ترتیب پیشنهادی می‌تواند «جای‌گذاری و مقیاس‌بندی عددها؛ جای‌گذاری و رمزگذاری گروه‌ها؛ اتصال خروجی‌ها؛ مدل» باشد. انتخاب ویژگی تحت نظارت نیز باید داخل همین فرایند ارزیابی قرار گیرد.

\section{رمزگذاری دسته‌ها، گسسته‌سازی و تاریخ}
برای ویژگی ترتیبی، کد عددی باید ترتیب را حفظ کند؛ بااین‌حال نسبت کد ۸ به ۴ الزاماً معنای «دو برابر بودن کیفیت» ندارد. برای رنگ‌های اسمی، عددهای ۱، ۲ و ۳ فاصله و ترتیب ساختگی می‌سازند. در رمزگذاری تک‌داغ \lr{One-Hot}، هر گروه یک نشانگر دارد. در مدل خطی با عرض از مبدأ، نگه داشتن همهٔ ستون‌های تک‌داغ می‌تواند وابستگی خطی ایجاد کند؛ حذف یک گروه مرجع یا منظم‌سازی، راه‌های معمول مدیریت آن‌اند. تکلیف گروه ناشناخته در آزمون باید مشخص باشد.

گسسته‌سازی، عدد را به بازه تبدیل می‌کند. در پهنای مساوی، طول بازه‌ها برابر است؛ در فراوانی مساوی، تعداد مشاهده‌ها تا حد امکان نزدیک است. تکرار مقدارها ممکن است مانع تقسیم دقیق به تعدادهای مساوی شود. مرزها نباید دو مقدار برابر را بی‌دلیل به گروه‌های متفاوت بفرستند. این تبدیل می‌تواند نویز را کم کند، ولی اطلاعات دقیق درون بازه را از بین می‌برد.

\begin{example}{دو روش بازه‌بندی}
برای دادهٔ مرتب $0,4,12,16,16,18,24,26,28$، سه بازهٔ با پهنای ۱۰، تعدادهای $2,4,3$ می‌دهند. با هدف سه مشاهده در هر گروه می‌توان مرزهای ۱۴ و ۲۱ را برگزید و سه گروه سه‌تایی ساخت. این مثال بدون شکستن تساوی مقدارهای ۱۶ امکان‌پذیر است، ولی چنین امکانی برای همهٔ داده‌ها وجود ندارد.
\end{example}

برای تاریخ، نخست قالب و منطقهٔ زمانی و سپس ویژگی کاربردی مانند روز هفته، فصل یا فاصله تا یک رویداد را استخراج کنید. ساعت ۲۳ و ۰ نزدیک‌اند؛ یک کد عددی خام این نزدیکی را نشان نمی‌دهد. نمایش چرخه‌ای با سینوس و کسینوس، برای دورهٔ $T$، می‌تواند چنین باشد:
\[\left(\sin\frac{2\pi t}{T},\cos\frac{2\pi t}{T}\right).\]
استفاده از متغیر زمانی باید با هدف پیش‌بینی سازگار باشد؛ اطلاعات آینده را نمی‌توان صرفاً با تغییر قالب، به ویژگی مجاز تبدیل کرد.

\section{از متن به بردار؛ فراوانی و TF-IDF}
در پاک‌سازی متن، نویسه‌ها یکدست و متن به واحدهای مناسب شکسته می‌شود. برای فارسی، تفاوت «ی/ي»، «ک/ك» و نیم‌فاصله مهم است. حذف نشانه‌گذاری یا واژهٔ توقف همیشه بی‌ضرر نیست: حذف «نه» می‌تواند معنی جمله را عوض کند. ریشه‌یابی شکل‌های صرفی را نزدیک می‌کند، ولی ممکن است واژه‌های متفاوت را نیز ادغام کند.

در کیسهٔ واژه‌ها \lr{Bag of Words}، ترتیب واژه‌ها نادیده گرفته می‌شود و هر سند با برداری از وقوع، شمارش یا فراوانی واژه‌ها نمایش داده می‌شود. نمایش دودویی تنها حضور را می‌سنجد؛ شمارش به طول سند حساس است. فراوانی نرمال‌شده و وزن‌دهی \lr{TF-IDF} اثر طول و واژه‌های بسیار رایج را تغییر می‌دهند:
\[w_{ij}=\operatorname{tf}_{ij}\log\frac{N}{\operatorname{df}_j}.\]
$N$ تعداد سندهای آموزش و $\operatorname{df}_j$ تعداد سندهای حاوی واژهٔ $j$ است. در قرارداد هموارشدهٔ رایج، وزن معکوس سند برابر است با
\[\operatorname{idf}_j=\log\frac{1+N}{1+\operatorname{df}_j}+1.\]
پس هنگام مقایسهٔ دو پیاده‌سازی باید قرارداد TF ، هموارسازی و نرمال‌سازی بردار را مشخص کرد؛ همهٔ خروجی‌های موسوم به TF-IDF عدد یکسانی نمی‌دهند.

\begin{example}{وزن واژهٔ عمومی و اختصاصی}
در سه سند، واژهٔ «داده» در هر سه و «خوشه» فقط در یکی آمده است. در قرارداد بدون هموارسازی، وزن IDF ِ «داده» صفر و وزن «خوشه» برابر $\log3$ است. در قرارداد هموارشده، وزن‌ها به‌ترتیب $1$ و $1+\log2$ می‌شوند. صفر نبودن واژهٔ عمومی در قرارداد دوم به معنی خطا نیست؛ تعریف وزن متفاوت است.
\end{example}

بردارهای واژه معمولاً تنک‌اند: از ده‌ها هزار مؤلفه تنها تعداد کمی ناصفر است. بردار تعبیه یا \lr{Embedding} نمایش چگال و آموخته‌شده‌ای است که رابطهٔ معنایی را به فاصله یا جهت بردار پیوند می‌دهد. مدل‌های رمزگذار مانند BERT ، رمزگشا مانند مدل‌های مولد و رمزگذارـ‌رمزگشا نقش‌های متفاوتی دارند؛ صرفاً کوتاه‌تر بودن بردار تضمین نمی‌کند که برای هر وظیفه بهتر باشد. انتخاب نمایش باید با داده و ارزیابی همان مسئله سنجیده شود.

\section{شباهت جاکارد، کسینوسی و معنای صفر مشترک}
برای بردار دودویی نامتقارن، نبودن هم‌زمان هزار واژه دلیل خوبی برای مشابه بودن دو سند نیست. جاکارد، حضورهای مشترک را نسبت به اجتماع حضورها می‌سنجد:
\[J(A,B)=\frac{|A\cap B|}{|A\cup B|}=\frac{M_{11}}{M_{11}+M_{10}+M_{01}}.\]
برای دو مجموعهٔ تهی، این کسر تعریف ندارد و قرارداد کاربرد باید جداگانه مشخص شود. ضریب تطابق ساده، صفرهای مشترک را هم می‌شمارد و برای دادهٔ تنک می‌تواند شباهت گمراه‌کننده بسازد.

برای بردارهای وزن‌دار غیرصفر، کسینوس زاویه به کار می‌رود:
\[\operatorname{cos}(\mathbf x,\mathbf z)=\frac{\mathbf x^\mathsf T\mathbf z}{\|\mathbf x\|_2\|\mathbf z\|_2}.\]
ضرب بردار در عدد مثبت، کسینوس را تغییر نمی‌دهد؛ جمع یک ثابت با همهٔ مؤلفه‌ها عموماً آن را تغییر می‌دهد. کسینوس برای بردارهای نامنفی بین صفر و یک است، ولی برای بردار عمومی می‌تواند منفی باشد. برای بردار صفر باید قاعدهٔ مشخصی در نظر گرفت.

\begin{example}{کنترل محاسبهٔ کسینوس}
برای $\mathbf x=(3,2,0,5,0,0,0,2,0,0)$ و $\mathbf z=(1,0,0,0,0,0,0,1,0,2)$، ضرب داخلی $5$ و مربع طول‌ها $42$ و $6$ است. بنابراین
\[\operatorname{cos}(\mathbf x,\mathbf z)=\frac5{\sqrt{42}\sqrt6}\approx0.315.\]
طول بردار دوم $\sqrt6\approx2.449$ است. محاسبهٔ طول‌ها و ضرب داخلی به‌صورت جدا، راه ساده‌ای برای پیدا کردن خطاهای عددی است.
\end{example}

\section{انتخاب ویژگی؛ جست‌وجوی یک زیرمجموعه}
ویژگی کم‌تغییر، نامرتبط یا تکراری ممکن است محاسبه را زیاد و تعمیم را ضعیف کند. حذف نام و شناسه یا ستون کاملاً ثابت، یک بررسی اولیه است؛ ولی کم بودن تغییر همیشه به معنی بی‌اهمیتی نیست، به‌ویژه برای رخداد نادر. همبستگی خطی نزدیک صفر نیز وابستگی غیرخطی را رد نمی‌کند.

روش فیلتر \lr{Filter} با معیاری مانند اطلاعات متقابل یا آزمون مناسب، ویژگی‌ها را مستقل از اجرای مکرر مدل هدف رتبه‌بندی می‌کند. روش پوششی \lr{Wrapper} زیرمجموعه را با اجرای مدل و سنجش اعتبارسنجی ارزیابی می‌کند. روش تعبیه‌شده \lr{Embedded} انتخاب را در خود آموزش انجام می‌دهد. فیلتر معمولاً ارزان‌تر است؛ پوششی به مدل و تقسیم داده حساس و پرهزینه‌تر است. ارزیابی همهٔ زیرمجموعه‌های $d$ ویژگی، $2^d$ حالت دارد.

در انتخاب پیش‌رونده از مجموعهٔ تهی شروع می‌کنیم و هر بار بهترین افزودن را می‌پذیریم. در حذف پس‌رونده از همهٔ ویژگی‌ها شروع می‌کنیم. این جست‌وجوهای حریصانه بهترین جواب جهانی را تضمین نمی‌کنند. برای مثال، در الگوی XOR هر ویژگی به‌تنهایی ممکن است بی‌اطلاع باشد، اما دو ویژگی با هم کاملاً مفید باشند؛ انتخاب پیش‌رونده‌ای که بدون بهبود تک‌ویژگی متوقف می‌شود، این تعامل را از دست می‌دهد. روش‌های حذف بازگشتی و انتخاب مبتنی بر اهمیت مدل نیز باید با اعتبارسنجی مناسب انجام شوند.

\section{تحلیل مؤلفه‌های اصلی و کاهش بُعد}
انتخاب ویژگی، ستون‌های موجود را برمی‌گزیند؛ \lr{PCA} ستون‌های تازه‌ای می‌سازد که ترکیب خطی ستون‌های قبلی‌اند. پس از مرکز کردن داده، مؤلفهٔ اول جهت واحدی است که واریانس تصویر داده روی آن بیشینه باشد:
\[\mathbf v_1=\underset{\|\mathbf v\|_2=1}{\operatorname{argmax}}\mathbf v^\mathsf T S\mathbf v,\qquad S=\frac{X_c^\mathsf T X_c}{n-1}.\]
مؤلفه‌های بعدی همین هدف را تحت قید عمود بودن بر قبلی‌ها دنبال می‌کنند. با تجزیهٔ ویژهٔ $S$، بردارهای ویژه جهت‌ها و مقدارهای ویژه واریانس‌ها را می‌دهند. برای $r$ مؤلفهٔ نخست:
\[Z=X_cV_r,\qquad \widehat X=ZV_r^\mathsf T+\boldsymbol\mu,\qquad \operatorname{EVR}_j=\frac{\lambda_j}{\sum_\ell\lambda_\ell}.\]
برگرداندن میانگین در رابطهٔ بازسازی، سطر به سطر انجام می‌شود. علامت هر بردار ویژه دلخواه است؛ برعکس شدن آن مفهوم مؤلفه یا بازسازی را تغییر نمی‌دهد. مؤلفه‌ها ناهمبسته‌اند، ولی بدون فرض اضافی الزاماً مستقل نیستند.

\begin{example}{PCA دو متغیر همبسته}
اگر $S=\left(\begin{smallmatrix}4&3\\3&4\end{smallmatrix}\right)$ باشد، مقدارهای ویژه $7$ و $1$ و جهت اول $(1,1)/\sqrt2$ است. یک مؤلفه $7/8=87.5\%$ واریانس را نگه می‌دارد. این جهت «بیشترین تغییر داده» را نشان می‌دهد، نه لزوماً «بهترین جهت برای پیش‌بینی برچسب». اگر برچسب در جهت کم‌واریانس جدا شود، حذف آن جهت می‌تواند به طبقه‌بندی آسیب بزند.
\end{example}

مقیاس ویژگی‌ها و داده‌های پرت بر PCA اثر دارند. انتخاب تعداد مؤلفه صرفاً با درصد واریانس، جای ارزیابی وظیفهٔ نهایی را نمی‌گیرد. ابتدا میانگین و جهت‌ها را از آموزش یاد بگیرید و همان‌ها را روی آزمون اعمال کنید.
\labref{در آزمایشگاه PCA ، جهت تصویر را بچرخانید و واریانس نگه‌داشته‌شده و خطای بازسازی را ببینید. سپس آن را با جهت مؤلفهٔ اول مقایسه کنید.}

\section{نمونه‌گیری و پوشش فضای داده}
کاهش تعداد سطرها با کاهش تعداد ستون‌ها متفاوت است. نمونهٔ تصادفی ساده احتمال انتخاب برابر می‌دهد؛ نمونهٔ طبقه‌بندی‌شده، سهم گروه‌ها را کنترل می‌کند. حفظ نسبت گروه‌ها برای برآورد توزیع مهم است، در حالی که برابر کردن تعداد گروه‌ها ممکن است برای آموزش یک مسئلهٔ نامتوازن هدف دیگری داشته باشد. «نمونهٔ متوازن» و «نمونهٔ نمایندهٔ توزیع» مترادف نیستند.

در روش \lr{Kennard--Stone}، ابتدا دورترین زوج انتخاب می‌شود؛ سپس در هر گام، برای هر نامزد فاصله‌اش تا نزدیک‌ترین عضو منتخب را می‌سنجیم و نامزدی با بزرگ‌ترین این فاصلهٔ کمینه را برمی‌گزینیم. این روش پوشش هندسی را تقویت می‌کند، ولی می‌تواند نقاط پرت را بیش از حد انتخاب و نسبت گروه‌ها را تغییر دهد. محاسبهٔ همهٔ فاصله‌های زوجی نیز هزینهٔ حافظه و زمان دارد.

نمونه‌گیری ماه دسامبر برای پیش‌بینی همهٔ سال یا نظرسنجی تلفنی تنها در ساعت کاری می‌تواند سوگیری انتخاب بسازد. هیچ الگوریتم پیچیده‌ای از داده‌ای که گروه مهمی را اصلاً نمی‌بیند، اطلاعات آن گروه را خلق نمی‌کند. منشأ نمونه و جمعیت هدف را همراه نتیجه گزارش کنید.
\section{اجرای یک خط لوله بدون نشت؛ گام به گام}
در چهار رکورد آموزش، سن‌ها $20,40,\text{missing},60$ هستند و سن یک مراجعه‌کنندهٔ تازه ۸۰ است. سه مقدار مشاهده‌شده میانگین ۴۰ دارند؛ مقدار گمشده را با همین میانگین جایگزین می‌کنیم. اکنون آموزش شامل $20,40,40,60$ است. اگر استانداردسازی با واریانسِ مخرج $n$ انجام شود:
\begin{align*}
\mu_{\mathrm{train}}&=40,\\
\sigma_{\mathrm{train}}&=\sqrt{\frac{(20-40)^2+0+0+(60-40)^2}{4}}=\sqrt{200},\\
z_{\mathrm{new}}&=\frac{80-40}{\sqrt{200}}\approx2.8284.
\end{align*}
ترتیب این خط لوله «جایگذاری، سپس مقیاس‌بندی» است؛ پارامترهای هر مرحله از خروجی مرحلهٔ قبل روی آموزش یاد گرفته می‌شوند. اگر فرد تازه را در میانگین و واریانس وارد کنیم، نمایشِ آموزش به پرسش وابسته می‌شود. در اعتبارسنجی چندبخشی، تمام این مراحل باید برای هر بخش از نو فقط با آموزش همان دور برازش شوند.

در تبدیل کمینه–بیشینه، آموزش کمینهٔ ۲۰ و بیشینهٔ ۶۰ دارد؛ برای فرد تازه مقدار $(80-20)/(60-20)=1.5$ به دست می‌آید. پس بازهٔ صفر تا یک، ویژگی دادهٔ آموزش است و تضمینی دربارهٔ هر دادهٔ آینده نیست. بریدن مقدار در یک، تصمیم جداگانه‌ای است که اطلاعات فاصله از محدودهٔ آموزش را از بین می‌برد. برای ستون ثابت هم مخرج صفر است و باید قرارداد، مثلاً تبدیل به صفر و ثبت ثابت بودن ستون، روشن باشد.

\subsection*{چرا کاهش بُعد همیشه بهبود پیش‌بینی نیست؟}
PCA جهت دارای بیشترین واریانس ورودی را پیدا می‌کند. فرض کنید یک ویژگی با پراکندگی زیاد تقریباً مستقل از هدف و ویژگی کوچک دیگری تعیین‌کنندهٔ هدف باشد. نگه داشتن فقط مؤلفهٔ پُرواریانس ممکن است اطلاعات پیش‌بینی را حذف کند. بنابراین تعداد مؤلفه‌ها را همراه با عملکرد خط لوله روی اعتبارسنجی بررسی می‌کنیم. واریانس توضیح‌داده‌شده، معیار بازسازی ورودی است؛ صحت طبقه‌بندی یا خطای هدف را مستقیماً اندازه نمی‌گیرد.

\section{خودآزمایی و تمرین}
\begin{selftests}
\begin{selftest}{تبدیلِ آموزش را روی مشاهدهٔ تازه اعمال کنید}
در مثال سن‌ها، مقدار استانداردشدهٔ ۸۰ و مقدار کمینه–بیشینهٔ آن چیست؟ آیا دومی خطاست؟
\answer{به‌ترتیب تقریباً $2.8284$ و $1.5$ است. دومی خطای محاسباتی نیست؛ مشاهده خارج از دامنهٔ آموزش است. برای محاسبهٔ آن نباید بیشینه را با استفاده از مشاهدهٔ تازه دوباره برازش کرد.}
\end{selftest}
\begin{selftest}{استاندارد با نرمال یکسان است؟}
آیا استانداردسازی یک توزیع دو‌قله‌ای را نرمال می‌کند؟
\answer{خیر؛ فقط مکان و مقیاس را تغییر می‌دهد. شکل چندقله‌ای باقی می‌ماند. میانگین صفر و واریانس یک، شرط کافی برای نرمال بودن نیست.}
\end{selftest}
\begin{selftest}{چرا واژگان و IDF هم از آموزش‌اند؟}
اگر پیش از تقسیم داده واژگان و وزن‌های TF-IDF را بسازیم، چه اطلاعاتی نشت می‌کند؟
\answer{حضور و فراوانی سندی واژه‌های آزمون در انتخاب نمایش اثر می‌گذارد. واژگان و IDF باید از اسناد آموزش یاد گرفته شوند؛ اسناد تازه با همان نگاشت تبدیل می‌شوند و واژهٔ ندیده قرارداد مشخصی دارد.}
\end{selftest}
\begin{selftest}{همبستگی و استقلال را جدا کنید}
اگر $X$ حول صفر متقارن و $Y=X^2$ باشد، آیا همبستگی صفر می‌تواند با وابستگی همراه باشد؟
\answer{بله؛ با وجود وابستگی تابعی کامل، تحت وجود گشتاورهای لازم $E[X^3]=0$ و کوواریانس صفر می‌شود. همبستگی خطی، همهٔ انواع وابستگی را نمی‌سنجد.}
\end{selftest}
\begin{selftest}{نمونهٔ متوازن یا نماینده؟}
در جامعه‌ای با یک درصد مثبت، برای آموزش تعداد دو طبقه را برابر می‌کنیم. آیا آزمون را هم باید برابر کنیم؟
\answer{اگر هدف، عملکرد در همان جامعه است، آزمون باید نمایندهٔ نرخ واقعی باشد. نمونه‌سازی آموزش تصمیمی جداست؛ متوازن کردن آزمون، درستی هشدار و برآورد فایدهٔ عملی را عوض می‌کند.}
\end{selftest}
\begin{selftest}{PCA و پوشش هندسی را تفسیر کنید}
آیا نگه داشتن ۹۵ درصد واریانس، حفظ ۹۵ درصد اطلاعات هدف را تضمین می‌کند؟ آیا Kennard--Stone الزاماً نسبت گروه‌ها را حفظ می‌کند؟
\answer{هیچ‌کدام تضمین نمی‌شود. PCA معمولی برچسب را نمی‌بیند و تنوع هندسی هم با نمایندگی توزیع یکی نیست؛ نقاط نادر یا پرت می‌توانند بیش از سهمشان انتخاب شوند. هر دو انتخاب را نسبت به هدف نهایی ارزیابی کنید.}
\end{selftest}
\end{selftests}
